\documentclass[10pt,itemsep]{beamer}
\usetheme{jambro}
\graphicspath{{./graphics/}}
\renewcommand{\memouser}{ACB}

% Authors/institutes/date/etc
%--------------------------------
\title[]{VAR Identification: A Practical Session}
\author[]{Ambrogio Cesa-Bianchi}
\date{\today}

\begin{document}

%=============================================================
\begin{frame}[plain]
    \titlepage{
        \begin{center}
            \begin{minipage}{0.8\textwidth}
            \end{minipage}
    \end{center}}
\end{frame}
%.......................................................
\shownotesblockzero{}

%=============================================================
\begin{frame}{Exercise}
    \small
    File \texttt{Exercise.m} loads the data and estimates a baseline VAR($12$) with four variables: CPI, IP, 1-year rate, and Excess Bond Premium (EBP) as in Gertler \& Karadi (2015) 
    
    \medskip

    \textbf{Objective}: Identify a monetary policy shock using three approaches and, for each,
    plot the impulse responses of all four variables:
    
    \medskip
    
    \begin{enumerate}
        \item \textbf{Cholesky} --- recursive identification 
        \item \textbf{Sign restrictions} 
        \item \textbf{External instruments} 
    \end{enumerate}
    
    \medskip
    
    \textbf{Compare} the three sets of impulse responses and discuss: which assumptions are needed for each approach? Which is most credible, and why? How do the results differ across approaches, and why?

\end{frame}
%.......................................................
\shownotesblock{
    Key points for discussion:
    (1) Cholesky identification hinges on the ordering; flipping Rate and EBP can
        change the results substantially --- illustrate the fragility.
    (2) Sign restrictions are set-identified: error bands tend to be wider.
    (3) IV avoids ordering and sign assumptions; its credibility rests on the
        instrument validity. Potential threat: ``information effects'' (Nakamura
        \& Steinsson 2018) --- the FF4 may capture Fed information, not just surprises.
    (4) GK prefer IV precisely because it is more agnostic about the structural shock
        configuration and is well-suited to trace the credit channel.
}

\end{document}
